\subsection{What Should the Model Be Trained Against?}
\label{section:fc_value_aware}

Throughout the preceding subsections the model has been trained against a likelihood or reconstruction loss. The Recurrent State-Space Model line uses a pixel reconstruction term, a reward log-likelihood, and a Kullback-Leibler consistency term between prior and posterior; the original Dyna model is a tabular maximum-likelihood estimate of the empirical transition; the controller-model architecture of \citet{HaSchmidhuber2018} fits its variational autoencoder and its mixture-density recurrent network to observation and next-latent likelihoods respectively. In each case the loss rewards the model for matching observed distributions, regardless of how those distributions enter the value calculation the agent actually optimizes. The objection raised in the present subsection is that the agent does not care about predicting the next observation perfectly; it cares about getting the value function right. This subsection traces the line that took that observation seriously, from value-aware model learning in \citet{FarahmandVAML2017} through the value-equivalence principle of \citet{GrimmVALEQ2020,GrimmPVE2021}, the value-targeted regret bound of \citet{AyoubVTR2020}, the MuZero instantiation of \citet{Schrittwieser2020MuZero}, and the recent calibration analysis of \citet{VoelckerCalib2025}.

\subsubsection{Value-aware model learning}
\label{section:fc_value_aware_vaml}

\citet{FarahmandVAML2017} formalize the objection as a loss. Fix a value function class $\mathcal{V}$, a data distribution $\mu$ over state-action pairs, and a model class within which the learner selects an estimate $\widehat P$ of the true transition kernel $P^\star$. The value-aware model learning loss is
\[ \mathcal{L}_{\text{VAML}}(\widehat P) = \sup_{V \in \mathcal{V}} \Big\Vert (\widehat P V)(\cdot) - (P^\star V)(\cdot) \Big\Vert_\mu, \]
where $(\widehat P V)(s, a) = \int V(s') \widehat P(\mathrm{d} s' \mid s, a)$ is the expected next-state value under the learned model and $(P^\star V)(s, a)$ is the same quantity under the true kernel. The loss penalizes the worst-case discrepancy in expected next-state value across $\mathcal{V}$, measured under $\mu$. A model that places mass on features of the next observation irrelevant to every value function in $\mathcal{V}$ is not penalized; a model that distorts the value expectation is. \citet{FarahmandIterVAML2018} embed the same idea inside approximate value iteration, replacing the supremum over a value class with the realized value iterate at each step, which yields the practical IterVAML estimator and a fitted error bound that propagates through the Bellman backups rather than through the model alone. \citet{AsadiWasserstein2018} read the construction through optimal transport, observing that under a Lipschitz assumption on the value class the VAML loss reduces to a Wasserstein distance between $\widehat P$ and $P^\star$, which connects the value-aware loss to the substantial empirical literature on Wasserstein generative modeling.

VAML is the formal statement of the intuition that the model should be wrong in ways that do not matter for the value functional. This is the technical answer to the conceptual openness of Sutton's original Dyna architecture, in which the model was permitted to be probabilistic and oftentimes incorrect (\S\ref{section:fc_dyna_q_architecture}) without specifying what kind of error was tolerable. The likelihood-trained Recurrent State-Space Model line (\S\ref{section:fc_rssm}) commits to a particular notion of error, namely the reconstruction loss on pixels and the Kullback-Leibler divergence on latents, and violates the value-aware principle whenever the pixel distribution is rich in detail orthogonal to the reward. Whether this matters in practice depends on the task; the empirical claim of the value-aware line is that it does, and that the gap widens as observation dimensionality grows.

\subsubsection{Value equivalence and MuZero}
\label{section:fc_value_aware_ve}

\citet{GrimmVALEQ2020} sharpen the VAML idea into a clean equivalence relation. Fix a policy set $\Pi$ and a value function set $\mathcal{V}$. Two transition kernels $\widehat P$ and $P^\star$ are value-equivalent with respect to $(\Pi, \mathcal{V})$ when their Bellman operators agree on every pair $(\pi, V) \in \Pi \times \mathcal{V}$, that is when $(T^\pi_{\widehat P} V)(s) = (T^\pi_{P^\star} V)(s)$ for all $s$ and all $(\pi, V)$ in the chosen sets. The equivalence class shrinks monotonically as $\Pi$ and $\mathcal{V}$ grow, collapsing to the singleton $\{P^\star\}$ in the limit. The practical implication is that a small model class can be sufficient for a restricted set of policies and value functions, even if it would be insufficient as a generative model of the environment. \citet{GrimmPVE2021} subsequently introduce proper value equivalence, which restricts the value class to fixed points of the chosen Bellman operators and gives the variant of value equivalence that aligns most directly with the MuZero training objective. \citet{AyoubVTR2020} push the value-aware idea into the regret-bound literature with value-targeted regression, which trains the model against value targets and matches the dimension-dependent rates of optimism-based reinforcement learning under standard assumptions.

The empirical instantiation is MuZero. \citet{Schrittwieser2020MuZero} train a latent recurrence jointly with a policy network and a value network against three losses applied at the next step of the recurrence, namely a reward prediction loss, a value prediction loss, and a policy prediction loss matched to the output of a Monte Carlo tree search rooted at the same latent. There is no reconstruction loss; the latent dynamics are never asked to decode observations. The transition kernel itself has no explicit target, since the only quantities that enter any loss are the reward, the value, and the policy at the rollout step. Empirically the construction matches AlphaZero on Go, chess, and shogi without access to the rules of those games, and reaches state-of-the-art performance on the Atari benchmark. \citet{Antonoglou2022StochMuZero} extend the construction to stochastic transitions by inserting a learned afterstate node, the game-tree analogue of a chance node, between the agent action and the next state. Read in economics terms, value-aware model learning treats the world model as a task-relevant misspecification, kept correct on the value functional even when it is wrong on the full observation distribution.\footnote{\citet{VoelckerCalib2025} revisit this family and observe that the standard sampled losses, including the canonical MuZero objective, are uncalibrated surrogate losses whose minima do not in general coincide with the minima of the underlying target loss on the value function; value-equivalence is a property of an in-distribution policy and value class, and a model trained purely against value targets has no incentive to remain accurate once the policy drifts. The paper proposes a calibrated variant that restores the correctness property in expectation and leaves the broader question of value-aware accuracy under policy drift open.}\footnote{The same move appears in the operations-research line on decision-focused learning, in which a predictor is trained with a loss aligned to the downstream decision quality rather than to a statistical scoring rule on the forecast itself \citep{DontiAmosKolter2017}. The inner solver there is a convex program rather than a Bellman backup, but the conceptual move, to align the loss with what the decision-maker actually needs and to treat likelihood as the wrong default whenever the model class is misspecified, is shared.}

\subsection{Convergence Point: TD-MPC2}
\label{section:fc_tdmpc2}

TD-MPC2 \citep{Hansen2024TDMPC2} is the convergence point of three lines that the preceding sections traced separately. Model predictive control with learned dynamics, in the style of PETS and MBPO (\S\ref{section:fc_mbpo}), supplies the planning frame. Recurrent latent world models, in the Dreamer line (\S\ref{section:fc_rssm}), supply the idea that the model should live in a learned latent and need not decode observations. Value-aware learning, articulated by MuZero and the value-equivalence literature (\S\ref{section:fc_value_aware}), supplies the target the model is trained against. TD-MPC2 keeps the planner of the first line, drops the observation decoder of the second, and adopts the value functional as the training signal of the third. The single empirical commitment that organizes the paper is that one fixed hyperparameter configuration should work across many tasks and many model sizes, with no per-task tuning and no rescaling of components as the network grows.

\subsubsection{Components and joint loss}
\label{section:fc_tdmpc2_components}

The agent learns five components jointly. An encoder $h$ maps an observation $s_t$ to a normalized latent $z_t = h(s_t)$. A latent dynamics model $d$ predicts the next latent, $\hat z_{t+1} = d(z_t, a_t)$. A reward model $R$ predicts the immediate reward from the latent and action, $\hat r_t = R(z_t, a_t)$. A terminal action-value function $Q$ supplies a bootstrap beyond the planning horizon, $\hat q_t = Q(z_t, a_t)$. A policy prior $p$ supplies an amortized action proposal $\hat a_t = p(z_t)$ used both to warm-start the planner and to seed the bootstrap targets. The joint loss is the sum of a latent consistency term that pushes the predicted next latent to agree with the encoded next observation, $\| d(z_t, a_t) - \operatorname{sg}(h(s_{t+1})) \|^2$ under a stop-gradient on the target encoding, a cross-entropy loss on the reward prediction, a cross-entropy loss on the temporal-difference targets for $Q$, and a behavioral cloning term that regresses $p$ on the actions selected by the planner. Notably absent is an observation reconstruction loss; the model is never asked to decode pixels or proprioceptive features and is trained only against quantities the value calculation actually consumes. This is the architectural commitment that distinguishes TD-MPC2 from the Dreamer line (\S\ref{section:fc_rssm}), in which the world model carries a reconstruction term throughout.

\subsubsection{MPPI planning with value bootstrap}
\label{section:fc_tdmpc2_mppi}

Action selection at each environment step uses Model Predictive Path Integral control, a sampling-based optimizer over Gaussian action sequences \citep{Hansen2024TDMPC2}. The planner maintains a sequence of Gaussian distributions $(\mathcal{N}(\mu_h, \sigma_h^2))_{h=t}^{t+H-1}$ over the next $H$ actions, draws a population of candidate sequences, rolls each one forward through the learned latent dynamics, scores the imagined returns, and refits $(\mu_h, \sigma_h)$ to a softmax-weighted average of the top-scoring sequences. The scoring rule is the finite-horizon return augmented with a terminal value bootstrap,
\[
  \mu^\star, \sigma^\star \;=\; \arg\max_{(\mu,\sigma)} \;\mathbb{E}_{(a_t, \ldots, a_{t+H}) \sim \mathcal{N}(\mu, \sigma^2)} \Big[\, \gamma^{H} Q(z_{t+H}, a_{t+H}) \;+\; \sum_{h=t}^{H-1} \gamma^{h-t} R(z_h, a_h) \,\Big],
\]
where the trajectory expectation is taken under the learned latent dynamics $d$ rolled from the current encoded state $z_t = h(s_t)$. The terminal $Q$ extends the planner's effective horizon past the finite window $H$, so the agent does not need to roll long open-loop sequences through the learned model to capture long-horizon credit. This is the architectural reason TD-MPC2 sidesteps the long-rollout compounding error that limits PETS (\S\ref{section:fc_mbpo}), which has no terminal value and must extend $H$ to capture distant reward. After the planner converges, the agent executes only $a_t \sim \mathcal{N}(\mu_t^\star, (\sigma_t^\star)^2)$ and re-plans at the next step. The amortized policy $p$ seeds the next planning call and supplies the action used in the $Q$-bootstrap.

\subsubsection{Empirical headline and scaling}
\label{section:fc_tdmpc2_empirics}

The paper evaluates TD-MPC2 across 104 continuous-control tasks spanning four benchmarks, DeepMind Control, Meta-World, ManiSkill2, and MyoSuite, under a single fixed hyperparameter configuration with no per-task tuning. The task suite includes action dimensionalities up to $39$, sparse rewards, multi-object manipulation, musculoskeletal motor control, and locomotion on dog and humanoid embodiments. Across all four benchmarks TD-MPC2 outperforms SAC, the original TD-MPC, and DreamerV3 at comparable parameter counts. The contribution is not that any single task is solved better than ever before, but that one configuration solves all of them, which is the property the field had previously assigned to model-free baselines under heavy per-task tuning.

The scaling experiment trains a single multi-task agent on $80$ tasks drawn from DeepMind Control and Meta-World and varies the world model size from $1$ million parameters to $317$ million. The normalized score moves from $16.0$ at $1$M parameters to $70.6$ at $317$M, roughly log-linear in parameter count over the two and a half orders of magnitude examined, with no observed saturation at the upper end. This is the first model-based reinforcement learning paper to report a clean monotone scaling curve over this range. The same hyperparameters are used at every model size; no learning rate, batch size, or capacity-specific schedule is retuned. The empirical claim of TD-MPC2 is that an MPC-shaped model-based agent, built from the components developed in the three earlier lines, scales with parameters in the same way that supervised and self-supervised models do.
